% Part: turing-machines
% Chapter: undecidability
% Section: unsolvability-decision-problem

\documentclass[../../../include/open-logic-section]{subfiles}

\begin{document}

\olfileid[tr]{tur}{und}{dec}
\olsection{Karar Problemi}

Verilen bir !!{sentence}{}nin geçerli olup olmadığını belirleyen etkili bir
yöntem varsa ve ancak varsa birinci dereceden mantığın \emph{karar
verilebilir} olduğunu söyleriz. Anlaşıldığı üzere böyle bir yöntem yoktur:
Birinci dereceden tümcelerin geçerliliğine karar verme problemi çözülemezdir.

Bu önemli olumsuz sonucu ortaya koymak için karar probleminin bir Turing
makinesiyle çözülemeyeceğini kanıtlarız. Yani şeridinde birinci dereceden bir
!!{sentence} bulunduğunda her zaman sonunda duran ve !!{sentence}{}nin geçerli
olup olmamasına göre $1$ ya da $0$ çıktısını veren hiçbir Turing makinesinin
bulunmadığını gösteririz. Church--Turing tezine göre hesaplanabilir her
fonksiyon Turing hesaplanabilirdir. Dolayısıyla bu ``geçerlilik fonksiyonu''
etkili biçimde hesaplanabilir olsaydı Turing hesaplanabilir olurdu. Turing
hesaplanabilir değilse etkili biçimde de hesaplanamaz.

Karar probleminin çözülemez olduğunu kanıtlama stratejimiz, durma problemini
ona indirgemektir. Bunun anlamı şudur: $e$ ile betimlenen Turing makinesi $w$
girdisinde duruyorsa $1$, aksi durumda $0$ çıktısıyla duran $h(e,w)$
fonksiyonunun Turing hesaplanabilir olmadığını kanıtladık. Birinci dereceden
tümcelerin geçerliliğine karar veren bir Turing makinesi bulunsaydı $h$'yi
hesaplayan bir Turing makinesinin de bulunacağını göstereceğiz. $h$ bir
Turing makinesince hesaplanamadığına göre geçerliliğe karar veren bir Turing
makinesi de bulunamaz.

Bu stratejinin ilk adımı, her $w$ girdisi ve $M$ Turing makinesi için $M$'nin
yönerge kümesiyle $w$ girdisini gösteren !!a{sentence} $!T(M,w)$ ile
``$M$ sonunda durur'' ifadesini dile getiren !!a{sentence}~$!E(M,w)$'yi
aşağıdaki koşulu sağlayacak biçimde etkili olarak betimleyebileceğimizi
göstermektir:
\begin{quote}
  $\Entails !T(M, w) \lif !E(M,w)$ ancak ve ancak $M$, $w$ girdisinde durursa.
\end{quote}
Kanıtımızın büyük bölümü bu $!T(M,w)$ ve $!E(M,w)$ tümcelerini betimlemekten
ve $!T(M,w)\lif !E(M,w)$ tümcesinin ancak ve ancak $M$, $w$ girdisinde
durursa geçerli olduğunu doğrulamaktan oluşacaktır.

\end{document}
